\bookchapter{第五十四章：項目一——我嘅記賬本（本地版）}

\chapterintro{歡迎嚟到你第一個「正式項目」。呢個記賬本雖然唔上網，但五臟俱全：記到一筆賬、睇到流水、按月統計到、又刪改到。做完佢，你就完整行咗一次「資料 邏輯 介面」嘅全流程。唔好心急抄程式碼——先諗清楚先至郁手，呢個係我哋一貫嘅原則。}

\booksection{一、需求思考：先問五條問題}

郁手之前，用「需求五要素」將個諗法捋清楚：

\begin{center}
\begin{tabularx}{\linewidth}{|X|X|}
\hline
\textbf{要素} & \textbf{答案} \tabularnewline
\hline
Who（畀邊個用） & 自己用，屋企人朋友都得 \tabularnewline
What（做乜嘢） & 記收支、睇流水、按月匯總、按分類統計 \tabularnewline
Why（點解） & 練手 + 搞清楚啲錢使咗去邊 \tabularnewline
How（點樣做） & Kotlin + Compose + Room + MVVM，純本地冇後端 \tabularnewline
Constraint（約束） & 一個禮拜之內完成；最低支援 Android 8；唔上架都先按上架標準寫 \tabularnewline
\hline
\end{tabularx}
\end{center}

\begin{center}\small\color{gray}記賬本需求五要素\end{center}

跟住拆解成「每一步都用得」嘅小步驟：(1) 建 Room 資料層（儲到一筆賬）；(2) 列表頁（睇到流水）；(3) 錄入頁（加到賬）；(4) 刪除同編輯；(5) 月度匯總卡。每完成一步，App 都係一樣「用得嘅嘢」。

\booksection{二、資料層：畀「賬」有個屋企}

先定義一筆賬係咩樣——欄位要夠用但唔好貪多：

\begin{lstlisting}[language=Kotlin]
@Entity(tableName = "records")
data class Record(
    @PrimaryKey(autoGenerate = true) val id: Long = 0,
    val type: String,        // "income" 收入 | "expense" 支出
    val amount: Double,      // 金額（元）
    val category: String,    // 分類：餐飲、交通、人工……
    val note: String = "",   // 備註（可選）
    val date: String         // "2026-08-15"
)
\end{lstlisting}

Dao 入面加兩條「統計用」嘅查詢——呢個係記賬本嘅靈魂：

\begin{lstlisting}[language=Kotlin]
@Dao
interface RecordDao {
    @Query("SELECT * FROM records ORDER BY date DESC, id DESC")
    fun getAll(): Flow<List<Record>>

    @Insert
    suspend fun insert(record: Record)

    @Delete
    suspend fun delete(record: Record)

    @Query("SELECT COALESCE(SUM(amount),0) FROM records WHERE type='income' AND date LIKE :month || '%'")
    suspend fun monthIncome(month: String): Double   // month 形如 "2026-08"

    @Query("SELECT COALESCE(SUM(amount),0) FROM records WHERE type='expense' AND date LIKE :month || '%'")
    suspend fun monthExpense(month: String): Double
}
\end{lstlisting}

\begin{tcbinfo}{小貼士：金額點解用 Double？}
嚴謹嘅金融 App 會用 BigDecimal 或者「以『分』為單位嘅 Int」嚟避免浮點誤差（0.1+0.2$\neq$0.3 呢個老問題）。練手項目用 Double 就夠；如果第日做真金融產品，記得轉做用「分」做單位——呢個係無數前輩踩出嚟嘅坑。

\end{tcbinfo}

\booksection{三、MVVM：介面同邏輯「分家」}

直接畀介面去操作資料庫？行到，但介面程式碼會又長又臭。MVVM 架構令大家各司其職：View（介面）淨係負責畫，ViewModel（邏輯中樞）管資料，Model（Room）管儲存。

View（Compose 界面）

ViewModel（狀態與邏輯）

Model（Room 數據層）

\begin{center}\small\color{gray}圖 5：MVVM 三層結構——資料單向流動，各層互唔過界\end{center}

\begin{lstlisting}[language=Kotlin]
class RecordViewModel(private val dao: RecordDao) : ViewModel() {

    // 將 Flow 轉成界面觀察到嘅狀態
    val records = dao.getAll()
        .stateIn(viewModelScope, SharingStarted.WhileSubscribed(), emptyList())

    // 用戶操作：調用 suspend 函數
    fun addRecord(type: String, amount: Double, category: String, note: String, date: String) {
        viewModelScope.launch {
            dao.insert(Record(type = type, amount = amount,
                category = category, note = note, date = date))
        }
    }

    fun deleteRecord(record: Record) {
        viewModelScope.launch { dao.delete(record) }
    }
}
\end{lstlisting}

介面（Compose）入面咁樣用：val vm: RecordViewModel = viewModel()，跟住 val list by vm.records.collectAsState()——資料一變，列表自動刷新，你乜都唔使做。呢個就係 MVVM + Flow + Compose 三劍合璧嘅威力。

\booksection{四、介面：列表頁 + 錄入頁}

列表頁用 LazyColumn（懶載入列表）——資料幾多都唔卡：

\begin{lstlisting}[language=Kotlin]
@Composable
fun RecordListScreen(vm: RecordViewModel) {
    val records by vm.records.collectAsState()

    Column {
        // 頂部彙總卡片（收入-支出=結餘）
        SummaryCard(records)

        // 流水列表
        LazyColumn {
            items(records, key = { it.id }) { record ->
                RecordRow(record) { vm.deleteRecord(record) }
            }
        }

        // 懸浮"記一筆"按鈕  彈錄入對話框
        FloatingActionButton(onClick = { showAddDialog = true }) {
            Text("+")
        }
    }
}

@Composable
fun RecordRow(record: Record, onDelete: () -> Unit) {
    Row(Modifier.fillMaxWidth().padding(12.dp)) {
        Text(record.category, fontWeight = FontWeight.Bold)
        Spacer(Modifier.weight(1f))
        Text(if (record.type == "income") "+${record.amount}" else "-${record.amount}",
             color = if (record.type == "income") Color(0xFF27AE60) else Color(0xFFC0392B))
    }
}
\end{lstlisting}

錄入用 AlertDialog（對話框）最慳事：幾個輸入框 + 收入／支出切換 + 分類下拉，撳確定就調 vm.addRecord(...)。

\begin{tcbwarn}{注意：輸入校驗唔好偷懶}
金額輸入框一定要校驗：空值、負數、「abc」呢類非法輸入全部要欄住（用 toDoubleOrNull() 判斷）。App 冧咗最常見嘅原因之一就係「將用戶嘅輸入當咗係乖仔」。記住：用戶輸入嘅所有嘢，都要當敵人咁審查。

\end{tcbwarn}

\booksection{五、編輯同刪除：畀資料「入得到又出得到」}

淨係識加唔識刪嘅記賬本，用三日就報廢。補返兩個操作：

\begin{lstlisting}[language=Kotlin]
// Dao 入面補兩條：
@Query("UPDATE records SET amount=:amount, category=:category, note=:note WHERE id=:id")
suspend fun update(id: Long, amount: Double, category: String, note: String)

@Query("DELETE FROM records WHERE id=:id")
suspend fun deleteById(id: Long)

// 界面上：長㩒列表項彈"編輯／刪除"菜單（Compose 用 DropdownMenu 或 AlertDialog）
// 刪除之前最好彈確認框——"刪錯嘢"係記賬類 App 用戶最憎嘅事
\end{lstlisting}

互動細節兩條經驗：刪除必須二次確認（防止手快快），編輯盡量預填返原本個值（改一個字比起重新輸入慳事）。唔好睇小呢兩個細節——產品感就係呢啲細節堆出嚟嘅。

\booksection{六、驗收：三層檢查 + 兩個進階挑戰}

\booksection{八、驗收：三層檢查 + 兩個進階挑戰}

項目做完，按「功能 邊界 架構」三層驗收：記唔記到、刪唔刪到、重開 App 資料仲喺唔喺度（功能）；金額輸 0 或者負數會唔會冧、轉屏會唔會冇咗資料（邊界）；介面同資料有冇撈埋一齊、欄位有冇寫死（架構）。

基礎版驗收過關之後，畀自己加兩味「進階菜」：

\begin{enumerate}
  \item 月度篩選：頂部加月份切換（例如 2026-08 v），列表淨係顯示當月流水；
  \item 分類圓餅圖：用 Compose 嘅 Canvas 或者引入輕量圖表庫，將「錢使咗去邊」畫成圓餅圖——做完呢樣嘢，記賬本就有啲「產品樣」喇。
\end{enumerate}

\begin{tcbpractice}{嚟動手：今個禮拜搞掂佢}
按「資料層 ViewModel 列表頁 錄入頁 匯總卡」嘅次序郁手，每日一步。卡咗就將報錯原封不動丟畀搜尋引擎或者 AI（仲記唔記得第七篇《AI編程實戰篇》？描述問題嗰陣講清楚「期望行為、實際行為、重現步驟」）。完成之後用它記一個禮拜賬——自己整嘅 App 用起上嚟，成就感會推住你繼續向前行。

\end{tcbpractice}

\begin{tcbsummary}{本章小結}
\begin{itemize}
  \item 需求五要素 拆成「每步用得」嘅小步驟，先至郁手
  \item 資料層：Record 實體 + Dao（包埋月份統計 SQL）+ Database
  \item MVVM 分工：View 淨係負責畫，ViewModel 管邏輯，Model 管儲存
  \item Flow + stateIn + collectAsState：資料一變介面自動刷新
  \item LazyColumn 懶載入列表；AlertDialog 做錄入
  \item 三層驗收（功能／邊界／架構）+ 進階挑戰（月度篩選、分類圓餅圖）
\end{itemize}

\end{tcbsummary}
